\section{量子回路の表現}
\label{sec:circuit-notation}

量子回路は，量子状態に加える演算と測定を回路図で表したものである．
本論文では量子線を横向きに描き，時間は左から右へ進むものとする．
1量子ビットゲートは量子線上の箱，制御ユニタリは制御点と標的の箱を結ぶ線，
測定はメータ形の記号で表す．複数量子ビットを一つのまとまりとして扱うとき，
この量子ビットのまとまりをレジスタという．

回路上で$U_1,U_2,\ldots,U_k$が左からこの順に現れる場合，
入力$\ket{\psi}$に対する出力は
\begin{equation}
 U_k\cdots U_2U_1\ket{\psi}
 \label{eq:circuit-composition-order}
\end{equation}
である．行列の積では，右端の演算から順に状態へ作用する．

\begin{prop}
ユニタリ回路$W=U_k\cdots U_1$の逆回路は
$W^\dagger=U_1^\dagger\cdots U_k^\dagger$である．
\end{prop}
\begin{proof}
行列の積の共役転置は$(AB)^\dagger=B^\dagger A^\dagger$であるから，
$(U_k\cdots U_1)^\dagger=U_1^\dagger\cdots U_k^\dagger$となる．
各$U_j$はユニタリなので$W^\dagger W=I$である．
したがって，逆回路を作るにはゲートの順を逆にし，
各ゲートをその共役転置に置き換えればよい．
\end{proof}

たとえば2量子ビットに$H\otimes I$を作用させた後，
第1量子ビットを制御とするCNOTを作用させる回路では，入力$\ket{00}$は
\[
 \ket{00}
 \xmapsto{H\otimes I}
 \frac{\ket{00}+\ket{10}}{\sqrt2}
 \xmapsto{\mathrm{CNOT}}
 \frac{\ket{00}+\ket{11}}{\sqrt2}
\]
と変化する．このように，本論文では重要な回路について，
回路図だけでなく各段階のケットも併記する．
